\documentclass{article}
\usepackage[margin=1in]{geometry}

\title{RotateAI Simulator Outline}
\author{Maxence Morel Dierckx}
\date{\today}

\begin{document}
\maketitle

\section{Note}
This document was the initial requirements outline. The implementation diverges in places. Please refer to the README and \verb|docs/derivations.pdf| for current behaviour.

\section{Problem Statement}
% What we need and why
RotateAI inference pipelines must run on STM32U5 microcontrollers for multi-month whale tag deployments at ultra low power. To design efficient programs, we need the ability to validate that a pipeline:

\begin{enumerate}
	\item Produces acceptable error in outputs.
	\item Meets resource constraints (space, time, power).
\end{enumerate}

The proposed simulator is a desktop tool that tests pipelines in the form of C binaries against real sensor data without requiring embedded hardware.

To enable rapid testing and iteration, the simulator must be decoupled from the pipelines, allowing any conforming binary to be tested with any input data.

\section{Functional Requirements}
% What the simulator must do

\subsection{Inputs}
% Loading prh.mat files, extracting A/M/p channels

The simulator will accept any quantity of MAT files with columns labelled:
\begin{itemize}
    \item $a_x$, $a_y$, $a_z$ (accelerometer).
    \item $m_x$, $m_y$, $m_z$ (magnetometer).
    \item $p$ (pressure).
    \item $aw_x$, $aw_y$, $aw_z$ (ground truth accelerometer).
    \item $mw_x$, $mw_y$, $mw_z$ (ground truth magnetometer).\\
\end{itemize}

\noindent The simulator will accept one reference to any compiled C binary.\\
	
\noindent The simulator will accept a configuration file in the TOML format providing:
\begin{itemize}
    \item Data sample rate in Hertz.
    \item Microcontroller voltage.
    \item Microcontroller Dhrystone Million Instructions Per Second (DMIPS) per MHz.
    \item Minimum frequency safety margin.
\end{itemize}

\subsection{Execution}
% How simulator calls the C pipeline, interface contract

The simulator will run a separate simulation for every MAT file and output a separate report for each.\\

\noindent A simulation will:
\begin{itemize}
    \item Run the binary.
    \item Feed data to the process sample by sample.
    \item Read data of the same shape from the process.
\end{itemize}

\subsection{Benchmarking}
% Time and space complexity measurement

The simulator will measure the following metrics:
\begin{itemize}
    \item \textbf{File size} of the program in bytes.
    \item \textbf{Peak memory usage} in bytes.
    \item \textbf{Average instructions} per inference.
    \item \textbf{Floating-point Operations} per inference.\\
\end{itemize}

\noindent The simulator will derive (estimate) the following metrics using configurable constants:
\begin{itemize}
    \item \textbf{Minimum operating frequency} in Hertz.
    \item \textbf{Energy per inference} estimate in Joules.
    \item \textbf{Duty cycle} as a percentage.
    \item \textbf{Power consumption} in Watts.\\
\end{itemize}
	
\noindent The simulator will validate the following metrics:
\begin{itemize}
    \item \textbf{Absolute error} from the ground truth in sensor units ($A_w$ in $g$, $M_w$ in $\mu T$).\\
\end{itemize}

\subsection{Outputs}
% Format of the output file

The simulator will produce one JSON file for every simulation.\\

\noindent Each JSON file will contain all of the stated metrics in human-readable format, with units labelled.\\

\section{Non-Functional Requirements}
% Performance, constraints, compatibility

The platform should be Linux only with Python 3.10+.\\

\noindent The simulation should take less than 100 times the real dataset duration to complete.\\

\noindent Identical inputs should produce identical outputs.\\

\section{Deliverable}
% Concrete outputs

Version 1.0.0 of the proposed simulator should be delivered as a repository on github.com/mtmd1 by 26/02/2026 at the latest.\\

\noindent It should include clear instructions in the form of a README documenting the entire installation process and intended workflow.\\

\noindent It should also contain a PDF document detailing how the derived metrics are estimated.\\

\section{Success Criteria}
% How we verify it works

The simulator will be considered successful when:

\begin{itemize}
    \item It runs without error on at least one provided test dataset.
    \item It produces a valid JSON file containing all specified metrics.
    \item Measured metrics (file size, instruction count) match independent verification
          (e.g., \verb|stat|, \verb|perf stat|).
    \item Derived metrics match hand calculations using the same input constants.
    \item A known-correct pipeline produces near-zero error; a known-incorrect pipeline produces significant error.
\end{itemize}

\end{document}